% ─────────────────────────────────────────────────────────────────────────────
% PISA Cognitive Tutor — Full Report
% Compiler: XeLaTeX (recommended for Thai/Chinese character support)
% Layout:   Two-column, A4, suitable for Overleaf
% ─────────────────────────────────────────────────────────────────────────────
\documentclass[twocolumn, 11pt, a4paper]{article}

% ── Fonts & Encoding (XeLaTeX) ────────────────────────────────────────────────

\usepackage{polyglossia}
\setmainlanguage{english}
% Uncomment the two lines below if compiling with pdfLaTeX instead:
% \usepackage[utf8]{inputenc}
% \usepackage[T1]{fontenc}

% ── Page Layout ───────────────────────────────────────────────────────────────
\usepackage[
    a4paper,
    top=2.2cm, bottom=2.2cm,
    left=1.5cm, right=1.5cm,
    columnsep=0.6cm,
    headheight=14pt
]{geometry}
\usepackage{fancyhdr}
\pagestyle{fancy}
\fancyhf{}
\fancyhead[L]{\small\textit{PISA Cognitive Tutor v2.0}}
\fancyhead[R]{\small\thepage}
\renewcommand{\headrulewidth}{0.4pt}

% ── Mathematics ───────────────────────────────────────────────────────────────
\usepackage{amsmath}
\usepackage{amssymb}
\usepackage{mathtools}

% ── Graphics & Figures ────────────────────────────────────────────────────────
\usepackage{graphicx}
\usepackage{float}
\usepackage{caption}
\usepackage{subcaption}
\captionsetup{font=small, labelfont=bf, skip=4pt}

% ── Tables ────────────────────────────────────────────────────────────────────
\usepackage{booktabs}
\usepackage{tabularx}
\usepackage{longtable}
\usepackage{array}
\usepackage{multirow}
\usepackage{makecell}
\renewcommand{\arraystretch}{1.2}

% ── Colors ────────────────────────────────────────────────────────────────────
\usepackage{xcolor}
\definecolor{codebg}{RGB}{248,248,248}
\definecolor{codegreen}{RGB}{0,128,0}
\definecolor{codepurple}{RGB}{128,0,128}
\definecolor{codeblue}{RGB}{0,0,180}
\definecolor{codegray}{RGB}{100,100,100}
\definecolor{placeholdercolor}{RGB}{220,230,245}
\definecolor{placeholderborder}{RGB}{100,130,180}

% ── Code Listings ─────────────────────────────────────────────────────────────
\usepackage{listings}
\lstdefinestyle{pythonstyle}{
    language=Python,
    backgroundcolor=\color{codebg},
    basicstyle=\scriptsize\ttfamily,
    keywordstyle=\color{codeblue}\bfseries,
    stringstyle=\color{codepurple},
    commentstyle=\color{codegreen}\itshape,
    numberstyle=\tiny\color{codegray},
    numbers=left, numbersep=5pt,
    breaklines=true, breakatwhitespace=false,
    frame=single, framerule=0.5pt,
    rulecolor=\color{gray!40},
    tabsize=4,
    showstringspaces=false,
    keepspaces=true,
    captionpos=b,
    extendedchars=true,
    literate={á}{{\'a}}1 {é}{{\'e}}1 {í}{{\'i}}1,  % extend as needed
}
\lstdefinestyle{sqlstyle}{
    language=SQL,
    backgroundcolor=\color{codebg},
    basicstyle=\scriptsize\ttfamily,
    keywordstyle=\color{codeblue}\bfseries,
    stringstyle=\color{codepurple},
    commentstyle=\color{codegreen}\itshape,
    numberstyle=\tiny\color{codegray},
    numbers=left, numbersep=5pt,
    breaklines=true,
    frame=single, framerule=0.5pt,
    rulecolor=\color{gray!40},
    tabsize=4,
    showstringspaces=false,
    captionpos=b,
}
\lstdefinestyle{textstyle}{
    backgroundcolor=\color{codebg},
    basicstyle=\scriptsize\ttfamily,
    breaklines=true,
    frame=single,
    tabsize=4,
    keepspaces=true,
    captionpos=b,
}
\lstset{style=pythonstyle}

% ── Hyperlinks ────────────────────────────────────────────────────────────────
\usepackage{hyperref}
\hypersetup{
    colorlinks=true,
    linkcolor=blue!60!black,
    citecolor=green!50!black,
    urlcolor=blue!70!black,
    pdftitle={PISA Cognitive Tutor v2.0},
    pdfauthor={Nutpraphut Prapatsirikul},
}

% ── Typography & Misc ─────────────────────────────────────────────────────────
\usepackage{microtype}
\usepackage{enumitem}
\setlist[itemize]{noitemsep, topsep=2pt}
\setlist[enumerate]{noitemsep, topsep=2pt}
\usepackage{titlesec}
\titleformat{\section}{\large\bfseries\scshape}{\thesection}{0.8em}{}[\titlerule]
\titleformat{\subsection}{\normalsize\bfseries}{\thesubsection}{0.7em}{}
\titleformat{\subsubsection}{\normalsize\itshape\bfseries}{\thesubsubsection}{0.6em}{}
\titlespacing{\section}{0pt}{8pt}{4pt}
\titlespacing{\subsection}{0pt}{6pt}{2pt}
\titlespacing{\subsubsection}{0pt}{4pt}{2pt}

% ── Image Placeholder Commands ────────────────────────────────────────────────
% Single-column placeholder
\newcommand{\imgplaceholder}[3][0.85\columnwidth]{%
    \begin{figure}[H]
        \centering
        \colorbox{placeholdercolor}{%
            \begin{minipage}{#1}
                \centering
                \color{placeholderborder}
                \vspace{1.2cm}
                {\large\textbf{[ IMAGE PLACEHOLDER ]}}\\[0.5em]
                \textit{#2}\\[1.2cm]
            \end{minipage}%
        }
        \caption{#2}
        \label{fig:#3}
    \end{figure}
}
% Double-column (wide) placeholder
\newcommand{\imgplaceholderwide}[2]{%
    \begin{figure*}[t]
        \centering
        \colorbox{placeholdercolor}{%
            \begin{minipage}{0.92\textwidth}
                \centering
                \color{placeholderborder}
                \vspace{1.8cm}
                {\large\textbf{[ IMAGE PLACEHOLDER ]}}\\[0.5em]
                \textit{#1}\\[1.8cm]
            \end{minipage}%
        }
        \caption{#1}
        \label{fig:#2}
    \end{figure*}
}

% ─────────────────────────────────────────────────────────────────────────────
% TITLE
% ─────────────────────────────────────────────────────────────────────────────
\title{%
    \LARGE\textbf{PISA Cognitive Tutor}\\[0.4em]
    \large A Multilingual AI-Powered Socratic Tutoring and\\
    Automated ICAP Cognitive Engagement Assessment System
}
\author{%
    \textbf{Nutpraphut Prapatsirikul}\\[0.3em]
    \small NECTEC (National Electronics and Computer Technology Center)\\[0.2em]
    \small \texttt{Version 2.0}
}
\date{May 22, 2026}

% ─────────────────────────────────────────────────────────────────────────────
\begin{document}

\maketitle
\thispagestyle{plain}

% ── Abstract ─────────────────────────────────────────────────────────────────
\begin{abstract}
This report presents the design, implementation, and evaluation of the PISA Cognitive
Tutor v2.0, a research-grade AI-powered Socratic tutoring system developed for NECTEC
to measure student cognitive engagement on PISA (Programme for International Student
Assessment) science questions. The system operationalizes the ICAP Framework
(Chi \& Wylie, 2014)---which categorizes learning modes as Passive, Active,
Constructive, and Interactive---into an automated, scalable assessment platform suitable
for publication-quality cognitive science research.

The architecture enforces a strict three-phase separation between (1)~offline document
ingestion, (2)~real-time Socratic tutoring, and (3)~offline automated evaluation. This
temporal decoupling prevents observer bias (Hawthorne Effect). A multilingual Vision
Language Model (Qwen3-VL-8B, 4-bit quantized) extracts structured knowledge from Thai,
English, and Chinese PISA PDFs. A hybrid retrieval engine combining dense and sparse
vector search with Reciprocal Rank Fusion over a Qdrant database, augmented by a
NetworkX knowledge graph, provides contextually rich scaffolding for the Socratic tutor.
A weighted ensemble of four locally-running LLMs (Qwen3, Mistral, Llama3.1, Typhoon2)
scores each student turn across eight cognitive dimensions (C1--C8) derived from the ICAP
taxonomy, with inter-rater agreement validated using three complementary statistics:
Fleiss' $\kappa$, Gwet's AC2, and Krippendorff's $\alpha$. Preliminary results indicate
that response latency is a credible proxy for cognitive depth (H1), that students tend
to progress toward higher ICAP levels over a session (H2), and that the ensemble
achieves Substantial agreement (AC2~$\geq$~0.60) on C1--C4 (H3). A FastAPI backend
exposes REST and WebSocket endpoints, and a Next.js~14 frontend provides the student
chat interface alongside a production-ready research dashboard with live cognitive
profiles, reliability statistics, and formula visualizations.
\end{abstract}

\tableofcontents
\vspace{1em}

% ─────────────────────────────────────────────────────────────────────────────
\section{Introduction}
% ─────────────────────────────────────────────────────────────────────────────

\subsection{Background Context}

The Programme for International Student Assessment (PISA), administered by the OECD,
is the world's most influential measure of fifteen-year-old students' applied literacy
in reading, mathematics, and science. Unlike curriculum-based tests, PISA emphasizes the
ability to use knowledge in real-world, open-ended contexts---precisely the kind of
higher-order thinking that predicts academic and professional success. PISA science items
require students to demonstrate competencies such as explaining phenomena scientifically,
evaluating and designing scientific inquiry, and interpreting data and evidence---competencies
that map naturally onto cognitive engagement frameworks from educational psychology.

In parallel, the fields of educational technology and learning analytics have seen a
proliferation of AI-powered tutoring systems. Intelligent Tutoring Systems (ITS) such as
Carnegie Learning's MATHia and Duolingo's adaptive language platforms use machine learning
to personalize content delivery. More recently, Large Language Models (LLMs) have enabled
conversational tutoring at scale. However, the vast majority of these systems optimize for
knowledge recall or skill mastery---measured by explicit test scores---rather than for the
underlying cognitive process: the depth and quality of a student's reasoning during the
tutoring interaction itself.

This project sits at the intersection of three converging trends: (1)~the global adoption
of PISA as a benchmark for educational quality, particularly in Southeast Asian nations
such as Thailand, (2)~the emergence of accessible, open-weight LLMs capable of running on
consumer-grade hardware, and (3)~growing research interest in real-time cognitive process
measurement as a complement to outcome-based assessment.

\subsection{Motivation}

The central motivation for PISA Cognitive Tutor v2.0 is the gap between what PISA measures
and how students are prepared. Traditional classroom instruction in Thailand and across the
ASEAN region remains predominantly passive: students receive information through lectures,
copy notes, and practice recall-oriented exercises. PISA items, by contrast, require
Constructive and Interactive cognitive engagement---the ability to construct explanations,
evaluate evidence, transfer knowledge across domains, and reflect on one's own reasoning.

There is currently no scalable infrastructure to (a)~deliver Socratic science tutoring
aligned to PISA competencies in a multilingual context (Thai, English, Chinese), (b)~automatically
measure the cognitive depth of student responses across a theoretically grounded taxonomy of
eight dimensions, and (c)~produce statistically validated inter-rater reliability data suitable
for inclusion in peer-reviewed cognitive science publications. Existing commercial systems are
closed, English-only, and do not expose the research-grade metrics needed to test specific
hypotheses about the relationship between engagement depth, response latency, and ICAP-level
progression. PISA Cognitive Tutor v2.0 was designed from first principles to fill exactly
this gap.

A secondary motivation is the research validation of a novel computational framework---the
Neuro-Software Engineering (C1--C8) approach---that maps each cognitive dimension to both a
neuroanatomical correlate and a corresponding behavior observable in student responses. This
framework, if validated, could serve as a bridge between cognitive neuroscience, educational
psychology, and software engineering education.

\subsection{Objectives and Scope}

\textbf{Primary objectives:}
\begin{enumerate}
    \item Build a Socratic tutoring system that scaffolds student reasoning on PISA science
    questions without revealing correct answers or providing evaluative feedback.
    \item Automatically score each student conversational turn across eight cognitive dimensions
    (C1--C8) using a validated ensemble of four LLMs, with scoring occurring strictly after
    the tutoring session ends to prevent observer effects.
    \item Produce statistically rigorous inter-rater reliability metrics (Fleiss' $\kappa$,
    Gwet's AC2, Krippendorff's $\alpha$) on the ensemble's scoring, targeting Substantial
    agreement (AC2~$\geq$~0.60) for the four foundational dimensions C1--C4.
    \item Test four research hypotheses (H1--H4) linking response latency to cognitive depth,
    turn-by-turn ICAP progression, ensemble reliability, and Thai-language model contribution.
    \item Expose all research metrics via a production-quality API and research dashboard.
\end{enumerate}

\textbf{Secondary objectives:}
\begin{enumerate}[resume]
    \item Extract structured PISA knowledge from multilingual PDFs using a Vision Language
    Model pipeline, producing a reusable knowledge base.
    \item Implement a hybrid retrieval system (dense + sparse vectors + knowledge graph) to
    provide PISA-aligned context to the Socratic tutor in real time.
    \item Deliver a production-quality Next.js research dashboard with student cognitive
    profiles, reliability tables, and KaTeX-rendered research formula visualizations.
\end{enumerate}

\textbf{Out of scope:}
\begin{itemize}
    \item Curriculum authoring tools or lesson planning features.
    \item Student authentication, authorization, or multi-tenancy.
    \item Real-time adaptive pedagogy based on individual student learning models.
    \item Validation of neuro-biomarker metrics (LCAI, CASM, MCI, PCRS) against physiological
    measurements.
    \item Deployment on cloud infrastructure; the system targets on-premise GPU servers.
\end{itemize}

\subsection{Document Structure}

Section~\ref{sec:challenge} articulates the three interlocking technical challenges that
shaped the architecture. Section~\ref{sec:background} provides the full technical
background: related literature, technology stack justification, and detailed methodology.
Section~\ref{sec:evaluation} presents the experimental setup, results, and critical
reflections. The Appendix contains a glossary, representative code snippets, extended
data tables, and the full bibliography.

% ─────────────────────────────────────────────────────────────────────────────
\section{Technical Challenge}
\label{sec:challenge}
% ─────────────────────────────────────────────────────────────────────────────

\subsection{Problem Formulation}

The core research problem can be stated as follows:

\begin{quote}
Given a sequence of student conversational turns $T = \{t_1, t_2, \ldots, t_n\}$ produced
during a Socratic tutoring session on a PISA science question, assign to each turn $t_i$ a
score vector $\mathbf{s}_i \in \{0, 1, 2, \text{NA}\}^8$, where each component corresponds
to one of eight cognitive dimensions C1--C8 derived from the ICAP Framework, such that the
ensemble scoring achieves Gwet's AC2~$\geq$~0.60 for dimensions C1--C4, without the scoring
process influencing the generation of $t_i$.
\end{quote}

This statement encodes four non-trivial sub-problems:

\textbf{Sub-problem 1: Multilingual Knowledge Extraction.}
PISA materials used by NECTEC exist as scanned or digitally typeset PDFs in Thai, English,
and Chinese. Standard PDF text extraction (PyMuPDF's text mode, pdfplumber) fails on Thai
scripts, mathematical notation, and tabular data. The problem is thus: how to extract
structured scientific knowledge from heterogeneous multilingual documents at research-grade
quality?

\textbf{Sub-problem 2: Contextual Socratic Generation.}
The tutoring LLM must produce responses that (a)~are grounded in PISA-specific scientific
knowledge, (b)~never directly reveal correct answers or use evaluative language, and
(c)~progressively deepen cognitive engagement toward higher ICAP levels.

\textbf{Sub-problem 3: Observer-Bias-Free Evaluation.}
If the scoring system runs during the tutoring session, the student's conversational
trajectory may be influenced, violating the ecological validity of the cognitive process
data. The evaluation must therefore be architecturally impossible to access from the tutoring
layer at runtime.

\textbf{Sub-problem 4: Reliable Ordinal Multi-Dimensional Scoring.}
Each of the eight cognitive dimensions must be scored on a three-point ordinal scale
(0/1/2) plus NA. A single LLM scorer is unreliable due to sensitivity to prompt phrasing
and stochastic sampling.

\subsection{Constraints and Trade-offs}

\textbf{Computational constraints.}
The system targets a single RTX~4090 (24~GB VRAM) research server. This imposes a hard
ceiling on total concurrent VRAM usage, driving the choice of 4-bit quantization for
Qwen3-VL-8B (reducing VRAM from $\approx$16~GB to $\approx$5~GB) and routing all LLM
inference through Ollama.

\textbf{Latency constraints.}
Socratic tutoring requires conversational latency under 10~seconds. Hybrid vector search
must complete in under 100~ms. LLM generation (streaming) must begin within 2~seconds of
receiving the user message. Evaluation is offline, so 10--30~seconds per turn is acceptable.

\textbf{Methodological constraint.}
The single hardest constraint is strict temporal and code-level separation between tutoring
and evaluation. The system has two completely independent data access paths---one for the
real-time chatbot (read-only from Qdrant, write-only to SQLite) and one for the offline
evaluator (read/write to SQLite scores table). Zero shared Python modules between these paths.

\textbf{Language constraints.}
Thai and Chinese PISA materials require multilingual handling throughout: the VLM must be
natively multilingual; the embedding model must produce semantically meaningful vectors for
Thai text; forbidden-word sanitization must cover three languages; and the evaluation
ensemble must include at least one Thai-tuned model (Typhoon2).

\textbf{Trade-off: accuracy vs.\ speed in ensemble voting.}
Running four 7--8B models in parallel takes 10--30~seconds. A single larger model
(e.g., Llama3-70B) takes 60--120~seconds but may achieve higher individual accuracy.
The four-model ensemble was chosen because model diversity reduces correlated errors and
$\kappa$-weighted voting is interpretable in the inter-rater reliability literature.

\subsection{Significance of the Challenge}

\textbf{Scientific contribution.} If the ensemble achieves AC2~$\geq$~0.60 on C1--C4
without human labeling, it validates LLM ensembles as a practical alternative to expensive
human annotation panels for cognitive process coding.

\textbf{Engineering contribution.} Strict architectural decoupling of tutoring and
evaluation, enforced at the Python import level, is a novel pattern for educational AI
systems---demonstrating that research validity constraints can be implemented as hard
technical invariants.

\textbf{Practical contribution.} A working multilingual PISA tutoring system that runs
entirely on-premise on a single GPU server, with no recurring API costs, is immediately
deployable in under-resourced educational research contexts across Southeast Asia.

% ─────────────────────────────────────────────────────────────────────────────
\section{Background, Technology, and Methodology}
\label{sec:background}
% ─────────────────────────────────────────────────────────────────────────────

\subsection{Literature Review}

\subsubsection{Historical Context and Prior Work}

\textbf{The ICAP Framework.} Chi, M.T.H., \& Wylie, R. (2014) propose that learning
modes exist on a hierarchy---Passive (receiving information), Active (manipulating
information), Constructive (generating new information), and Interactive (co-construction
in dialogue)---with each level producing better learning outcomes than the one
below~\cite{chi2014icap}. The framework provides the four-category taxonomy that the
C1--C8 scoring rubric operationalizes.

\textbf{Automatic Dialogue Act Classification.} Shriberg et al.\ (2004) established that
human rater agreement on dialogue categories rarely exceeds $\kappa = 0.70$~\cite{shriberg2004},
providing a reference baseline for automated agreement. This informed the choice of Gwet's
AC2 as the primary metric (robust to the Kappa Paradox in skewed distributions).

\textbf{Socratic Tutoring Systems.} Graesser et al.'s AutoTutor~\cite{graesser2004}
demonstrated that conversational tutoring systems could improve learning outcomes relative
to reading alone. However, AutoTutor is English-only, uses pre-scripted dialogue templates,
and was developed before the LLM era.

\textbf{Hybrid Dense-Sparse Retrieval.} Lin et al.\ (2021)~\cite{lin2021} establish that
combining dense neural embeddings with sparse keyword-based scoring consistently outperforms
either approach alone. The Reciprocal Rank Fusion (RRF) algorithm~\cite{cormack2009} used
in this system requires no learned fusion weight and remains stable across query distributions.

\textbf{Multilingual Embeddings.} Chen et al.\ (2024)~\cite{chen2024bge} present BGE-M3,
the state-of-the-art multilingual embedding model used in this system. It produces both
dense (1024-dimensional) and sparse (SPLADE-style) embeddings from a single forward pass,
natively supporting 100+ languages including Thai.

\textbf{Vision Language Models for Document Understanding.} Wang et al.\ (2024)~\cite{wang2024qwen2vl}
demonstrate the ability to process high-resolution document images with accurate extraction
of structured text, tables, and diagrams across multiple languages, avoiding the OCR error
cascade of traditional PDF extraction pipelines.

\subsubsection{Limitations of Existing Solutions}

\textbf{AutoTutor and rule-based ITSs} rely on manually authored dialogue scripts and
predefined knowledge bases. They cannot handle open-ended PISA science questions, do not
support non-English languages, and produce no cognitive process data.

\textbf{Commercial LLM tutors} (e.g., Khan Academy's Khanmigo, Duolingo Max) are
English-only, operate on closed infrastructure, and do not expose raw cognitive process data
or inter-rater reliability statistics.

\textbf{GPT-4/Claude-based single-model scorers} produce no natural inter-rater agreement
estimate, are sensitive to prompt phrasing, and have no mechanism for uncertainty
quantification. The four-model ensemble with $\kappa$-weighting directly addresses this gap.

\textbf{Closed-source educational analytics platforms} provide engagement metrics (login
frequency, video watch time, quiz scores) but cannot measure the cognitive depth of
individual written responses.

\subsection{Technology Stack}

\subsubsection{Core Frameworks and Languages}

\begin{table}[H]
\centering
\caption{Technology Stack by Layer}
\label{tab:tech-stack}
\scriptsize
\begin{tabularx}{\columnwidth}{@{}lXl@{}}
\toprule
\textbf{Layer} & \textbf{Technology} & \textbf{Version} \\
\midrule
Primary language & Python & 3.11 \\
Frontend language & TypeScript & 5.x \\
Vision LLM (ingestion) & Qwen3-VL-8B-Instruct & Latest \\
LLM inference runtime & Ollama & 0.3+ \\
Chat LLM & Qwen3-8B-Instruct & Latest \\
Embedding model & BAAI/bge-m3 & 1.2.9 \\
Vector database & Qdrant & 1.9.2 \\
Knowledge graph & NetworkX DiGraph & 3.3 \\
Session database & SQLite (WAL mode) & 3.x \\
Async DB driver & aiosqlite & 0.20 \\
PDF rendering & PyMuPDF (fitz) & 1.24+ \\
Quantization & bitsandbytes (4-bit QLoRA) & Latest \\
Web API & FastAPI & 0.111+ \\
ASGI server & Uvicorn & 0.30+ \\
Frontend & Next.js 14 (App Router) & 14.x \\
Statistical analysis & SciPy, NumPy, scikit-learn & 1.13+, 1.26+, 1.4+ \\
Visualization & Matplotlib, Seaborn & 3.9+, 0.13+ \\
\bottomrule
\end{tabularx}
\end{table}

\subsubsection{Justification for Technological Choices}

\textbf{Qwen3-VL-8B over alternatives (GPT-4V, LLaVA-1.6, Idefics2).}
GPT-4V would introduce recurring per-page API costs prohibitive for a multi-hundred-page
PISA corpus and violates the on-premise requirement. Qwen3-VL-8B runs entirely locally on
the RTX~4090, supports Thai/Chinese/English natively, and achieves state-of-the-art
performance on multi-page document understanding benchmarks.

\textbf{BGE-M3 over text-embedding-3-large or multilingual-e5.}
text-embedding-3-large is English-centric and requires an OpenAI API call per record.
multilingual-e5 produces only dense embeddings. BGE-M3 uniquely produces both dense and
sparse embeddings from one forward pass, making it the only viable choice for native hybrid
retrieval on Thai/English/Chinese text.

\textbf{Qdrant over Weaviate, Pinecone, or pgvector.}
Qdrant supports true hybrid search (dense + sparse vectors with native RRF fusion) as of
version 1.9. Qdrant runs locally without a Docker dependency; Pinecone is hosted-only,
violating the on-premise requirement.

\textbf{SQLite over PostgreSQL.}
The session database is append-dominant, single-writer, and read by at most a few
researchers at a time. SQLite in WAL mode handles this perfectly with zero operational
overhead.

\textbf{Four-model ensemble over a single larger model.}
Diverse 7--8B models produce uncorrelated errors, and when three of four models agree,
majority vote is more reliable than any single model. The $\kappa$-weighting mechanism
amplifies the signal from historically better-agreeing models. Four parallel 7--8B models
complete evaluation in 10--30~seconds vs.\ 60--120~seconds for a single 70B model.

\subsection{Methodology Overview}

\subsubsection{System Architecture}

The system enforces a strict three-phase architecture with no cross-phase imports at the
Python module level. Figure~\ref{fig:architecture} illustrates the full pipeline.

\imgplaceholderwide{Three-Phase System Architecture: Ingestion \textrightarrow{} Tutoring \textrightarrow{} Evaluation (replace with rendered pipeline diagram)}{architecture}

\textbf{Phase 1 --- Ingestion (offline, runs once or incrementally).}
PISA PDFs are rendered to page images at 200~DPI by PyMuPDF and grouped into 3-page
shards. Each shard is processed by Qwen3-VL-8B (4-bit QLoRA) to extract structured JSON
(topic, competency, misconceptions, vocabulary). Records are appended to
\texttt{knowledge\_base.jsonl}; BGE-M3 encodes each record to dual vectors that are
upserted to Qdrant and serialized to a NetworkX DiGraph (\texttt{knowledge\_graph.gpickle}).

\textbf{Phase 2 --- Tutoring (real-time, per-session).}
Each student turn triggers: (1)~duration measurement, (2)~Qdrant hybrid search (top $k=5$),
(3)~knowledge graph subgraph retrieval (2-hop BFS, up to 20 nodes), (4)~system prompt
assembly, (5)~streaming inference from Ollama (qwen3:8b), (6)~multilingual forbidden-word
sanitization, and (7)~turn persistence to SQLite. Sessions close after 12~turns.

\textbf{Phase 3 --- Evaluation (offline, after sessions complete).}
The \texttt{AnalyzerPipeline} fetches unscored turns via \texttt{LEFT JOIN}, retrieves
PISA context from Qdrant, and passes each turn to the \texttt{MajorityVoter} which runs
four LLMs in parallel (\texttt{ThreadPoolExecutor}). Kappa-weighted majority voting
produces scores for each C1--C8 dimension. Scores are saved to the SQLite scores table.
The metrics layer then computes $\kappa$, AC2, $\alpha$, LTA matrices, and H1--H4
statistics.

\subsubsection{Algorithm Design}

\textbf{Reciprocal Rank Fusion (RRF) for hybrid retrieval.}
Given a query $q$, the system produces two ranked lists from Qdrant: $L^d$ (dense/semantic)
and $L^s$ (sparse/keyword), each of length $2k$. The fused score for a document $d$ is:

\begin{equation}
    \text{RRF}(d) = \sum_{r \in \{\text{dense},\,\text{sparse}\}} \frac{1}{\text{rank}_r(d) + 60}
    \label{eq:rrf}
\end{equation}

The constant 60 is the standard smoothing factor from Cormack et al.~(2009)~\cite{cormack2009}.

\imgplaceholder{Hybrid Retrieval Pipeline: Dense + Sparse Vectors fused via RRF (replace with system diagram)}{retrieval}

\textbf{Kappa-weighted majority voting.}
For a given turn $t$ and dimension $C_j$, each model $m_i$ provides a vote
$v_i \in \{0, 1, 2, \text{NA}\}$ and a confidence $c_i \in [0, 1]$. Model weight:

\begin{equation}
    w_i = \kappa_i \times c_i
    \label{eq:kappa-weight}
\end{equation}

where $\kappa_i$ is the model's historical Gwet's AC2 agreement with the ensemble majority
on past turns (initialized to 1.0, then updated incrementally). Weighted vote totals:

\begin{equation}
    W_v = \sum_{\{i:\, v_i = v\}} w_i, \quad v \in \{0, 1, 2\}
\end{equation}

The dominant score is $\operatorname{argmax}_v W_v$. A turn is marked \emph{controversial}
if $\max_v W_v \,/\, \sum_v W_v < 0.50$.

\textbf{Engagement trajectory detection.}
Let $T_1 = \text{first } \lfloor n/2 \rfloor$ turns and $T_2 = \text{last } \lfloor n/2 \rfloor$
turns. High-order engagement score for a half:

\begin{equation}
    E(T) = \underset{t \in T}{\text{mean}} \; \underset{j=4}{\overset{8}{\text{mean}}} \; \text{score}(t, C_j)
\end{equation}

Trajectory is \emph{improving} if $E(T_2) - E(T_1) > 0.2$, \emph{declining} if
$E(T_1) - E(T_2) > 0.2$, and \emph{stable} otherwise.

\subsection{Database Architecture}

\subsubsection{Data Modeling and Schema}

The SQLite database at \texttt{data/chat\_logs/sessions.db} contains three tables.

\textbf{\texttt{sessions} table} --- one row per tutoring session:

\begin{lstlisting}[style=sqlstyle, caption={Sessions table schema}]
CREATE TABLE sessions (
    session_id  TEXT PRIMARY KEY, -- UUID v4
    user_id     TEXT NOT NULL,    -- UUID v4 (de-identified)
    started_at  TEXT NOT NULL,    -- ISO 8601 UTC
    question_id TEXT              -- optional PISA question ID
);
\end{lstlisting}

\textbf{\texttt{turns} table} --- one row per conversational exchange:

\begin{lstlisting}[style=sqlstyle, caption={Turns table schema}]
CREATE TABLE turns (
    id             INTEGER PRIMARY KEY AUTOINCREMENT,
    session_id     TEXT NOT NULL REFERENCES sessions,
    user_id        TEXT NOT NULL,
    turn_number    INTEGER NOT NULL,
    timestamp      TEXT NOT NULL,
    user_input     TEXT NOT NULL,
    ai_response    TEXT NOT NULL,
    kb_context_used INTEGER NOT NULL,
    duration_sec   REAL,
    question_topic TEXT
);
CREATE INDEX idx_turns_session ON turns(session_id);
CREATE INDEX idx_turns_user    ON turns(user_id);
\end{lstlisting}

\textbf{\texttt{scores} table} --- one row per scored turn (key columns shown):

\begin{lstlisting}[style=sqlstyle, caption={Scores table schema (abbreviated)}]
CREATE TABLE scores (
    id             INTEGER PRIMARY KEY AUTOINCREMENT,
    turn_id        INTEGER NOT NULL REFERENCES turns(id),
    session_id     TEXT NOT NULL,
    dominant_icap  TEXT, -- 'Passive'|'Active'|...|'Interactive'
    models_responded INTEGER,
    elapsed_sec    REAL,
    scored_at      TEXT NOT NULL,
    -- 8 dimensions x 3 cols = 24 score columns:
    C1_dominant TEXT, C1_max TEXT, C1_controversial INTEGER,
    -- ... C2 through C8 analogously ...
    C8_dominant TEXT, C8_max TEXT, C8_controversial INTEGER,
    individual_votes_json TEXT  -- full ensemble JSON
);
\end{lstlisting}

\textbf{Normalization.} The schema is in Third Normal Form (3NF). \texttt{user\_id} and
\texttt{turn\_number} are stored redundantly in \texttt{scores} as a deliberate
denormalization for analytical query convenience. A parallel \texttt{conversations.csv}
flat-file export is maintained for import into R, Stata, or SPSS.

\subsubsection{Data Ingestion and Pipeline}

Each validated JSON record from \texttt{knowledge\_base.jsonl} is canonicalized as:

\begin{lstlisting}[style=textstyle, caption={Canonical text string format}]
{topic} | {question_summary} | {correct_concept}
       | {common_misconceptions} | {key_vocabulary}
\end{lstlisting}

BGE-M3 encodes this in batches of 32, producing a 1024-dimensional dense vector and a
SPLADE sparse vector. Each record is upserted to Qdrant as a \texttt{PointStruct} with
\texttt{vector["dense"]} (float32 array) and \texttt{vector["sparse"]}
(\texttt{SparseVector}). Qdrant automatically builds an HNSW index (dense ANN) and an
inverted index (sparse matching) stored at \texttt{data/vector\_db/qdrant/}.

\textbf{Resumable ingestion.} The orchestrator reads existing \texttt{knowledge\_base.jsonl}
on startup and extracts already-processed \texttt{source\_file} values. Any PDF matching
this set is skipped, allowing safe interruption and resumption of multi-hour ingestion jobs.

\subsubsection{Security and Governance}

User IDs are UUID~v4 values generated fresh per session with no link to real names or
institutional identifiers. No PII is stored. All writes use parameterized queries throughout
the codebase---never string concatenation---eliminating SQL injection risk. \texttt{aiosqlite}
provides async, non-blocking writes during the tutoring phase. Foreign key constraints are
enforced (\texttt{PRAGMA foreign\_keys = ON}).

\subsection{Chatbot Architecture and Integration}

\subsubsection{Natural Language Processing Pipeline}

The chatbot uses a retrieval-augmented generation (RAG) approach:

\begin{enumerate}
    \item \textbf{Hybrid search} retrieves the top-5 most relevant PISA knowledge base
    records (semantic + keyword via BGE-M3 + Qdrant).
    \item \textbf{Knowledge graph enrichment} runs a 2-hop BFS from the best-matching
    concept node, returning up to 20 related nodes and their relationships
    (e.g., \texttt{HAS\_MISCONCEPTION}, \texttt{TESTS\_COMPETENCY}).
    \item \textbf{Context assembly} formats retrieved records and appends the graph
    subgraph summary.
    \item \textbf{Prompt construction} injects context into the system prompt alongside
    Socratic scaffolding instructions, the forbidden word list, and the last 12 conversation
    turns.
    \item \textbf{LLM generation} (Ollama/Qwen3-8B) produces a 2--4 sentence Socratic
    response ending in a single guiding question.
    \item \textbf{Post-generation sanitization} applies a multilingual regex over the
    output and replaces any match against forbidden evaluative words with
    \texttt{"..."}, logging the event.
\end{enumerate}

Prompt engineering parameters: \texttt{temperature=0.7}, \texttt{top\_p=0.9},
\texttt{max\_tokens=1024}. The system prompt enforces a ``one question per turn''
constraint to prevent cognitive overload.

\subsubsection{Dialogue Management and State Tracking}

Conversation history is maintained in two places:

\textbf{In-memory (per request).} The last 12 turns are held as a
\texttt{List[Dict]} with \texttt{role} and \texttt{content} keys, passed directly to the
LLM as the \texttt{messages} array.

\textbf{Persistent (SQLite).} Each turn is written to the \texttt{turns} table after the
AI response completes. On session resume, history is reconstructed from the database.

\textbf{12-turn maximum.} Research on PISA-aligned Socratic tutoring sessions indicates
that cognitive engagement metrics stabilize by turn~8--10. Twelve turns was chosen to
capture the full engagement trajectory while preventing fatigue effects.

\subsubsection{System Integration}

\textbf{FastAPI backend} (\texttt{src/api/main.py}) runs on port~8001, exposing REST
endpoints for session management, file upload, history retrieval, and profile queries,
plus a \texttt{/api/chat/\{session\_id\}} WebSocket endpoint that streams tokens
incrementally. CORS is configured for \texttt{http://localhost:3000}.

\textbf{Next.js~14 frontend} (\texttt{frontend/}) runs on port~3000 and serves both
the student chat interface and the researcher dashboard, communicating with the FastAPI
backend over HTTP and WebSocket.

\imgplaceholder{Next.js Chat Interface --- Socratic tutoring session in progress (replace with actual screenshot)}{nextjs-ui}

\imgplaceholder{Next.js Research Dashboard --- Live cognitive profiles, reliability statistics, and formula visualizations (replace with actual screenshot)}{nextjs-dashboard}

% ─────────────────────────────────────────────────────────────────────────────
\section{Evaluation and Discussion}
\label{sec:evaluation}
% ─────────────────────────────────────────────────────────────────────────────

\subsection{Experiment}

\subsubsection{Experimental Setup}

\textbf{Hardware environment:}
\begin{itemize}
    \item GPU: NVIDIA RTX~4090 (24~GB GDDR6X VRAM)
    \item CPU: AMD Ryzen~9 7950X (16 cores / 32 threads)
    \item RAM: 64~GB DDR5-6000
    \item Storage: NVMe SSD, 2~TB
    \item OS: Ubuntu~24.04 (production)
\end{itemize}

\textbf{Dataset.} PISA science questions provided by NECTEC covering three thematic domains:
Life Science (photosynthesis, ecosystems, genetics), Physical Science (energy transfer,
forces, electromagnetic radiation), and Earth and Space Science (plate tectonics, climate
systems, stellar evolution). Total corpus: $\approx$120~PDF pages across 8~PDF files in
Thai, English, and Chinese. Ingestion produced $\approx$340 validated knowledge base records.

\textbf{Student population.} 24~students from a pilot cohort (Grade~10, age~15--16) at a
NECTEC partner school in Bangkok. Each student completed one session of up to 12~turns.
Total turns collected: 271 (some sessions ended early due to student inactivity).

\textbf{Gold standard.} A subset of 80~turns was independently coded by two domain-expert
human annotators (a science education researcher and a cognitive psychologist, both fluent
in Thai and English) using the same C1--C8 rubric provided to the LLM ensemble. This subset
was held out from the $\kappa$-weight adaptation phase to serve as an unbiased reliability
baseline.

\subsubsection{Evaluation Metrics}

\textbf{For H1 (duration-depth correlation):}
Spearman rank correlation $\rho(\overline{C4\text{--}C8},\; \text{duration\_sec})$ per turn
across all 271 turns. Significance threshold: $p < 0.05$. Target: $\rho \geq 0.35$.

\textbf{For H2 (ICAP progression):}
Latent Transition Analysis: compute 4$\times$4 empirical transition probability matrix
$P$ where $P[i,j] = \Pr(\text{ICAP}_{t+1} = j \mid \text{ICAP}_t = i)$. Positive diagonal
dominance: $P[i,i+1] > P[i,i-1]$ for $i \in \{\text{Passive, Active, Constructive}\}$.

\textbf{For H3 (ensemble reliability):}
Gwet's AC2 (primary), Fleiss' $\kappa$, and Krippendorff's $\alpha$ per dimension C1--C8
against the 80-turn human annotation subset.
Target: AC2~$\geq$~0.60 for C1--C4; AC2~$\geq$~0.50 for C5--C8.

\textbf{For H4 (Thai model contribution):}
Per-model agreement rate on turns where \texttt{question\_topic} was Thai-language.
Target: Typhoon2's agreement rate $\geq$ mean of other three models $+ 5$~pp.

\textbf{Additional metrics:}
Forbidden word slip rate, knowledge base utilization rate, average response latency, models
responded rate, average evaluation latency per turn, and controversy rate.

\subsubsection{Test Execution}

\textbf{Scenario~A --- Baseline calibration ($N=20$ turns).}
A convenience sample drawn uniformly across ICAP levels was used to establish initial
$\kappa$ weights. Human annotations were provided; model weights were updated based on
agreement with the human gold standard.

\textbf{Scenario~B --- Full pilot session evaluation ($N=251$ turns).}
The remaining 251~turns were scored by the ensemble using $\kappa$-adapted weights,
processed in two batches: batch~1 (turns~1--8 from all sessions) and batch~2 (turns~9--12).

\textbf{Scenario~C --- Stress test (10 concurrent sessions).}
Ten concurrent WebSocket chat sessions were simulated using \texttt{httpx.AsyncClient},
each sending a pre-scripted 6-turn conversation. Metrics: $p_{50}$, $p_{95}$, $p_{99}$
response latency and error rate.

\textbf{Scenario~D --- Forbidden word compliance.}
All 271~AI responses were retrospectively analyzed by the sanitizer regex to measure the
pre-sanitization slip rate.

\begin{table}[H]
\centering
\caption{Raw Experimental Results Summary}
\label{tab:results}
\scriptsize
\begin{tabularx}{\columnwidth}{@{}Xll@{}}
\toprule
\textbf{Metric} & \textbf{Value} & \textbf{Target} \\
\midrule
H1: Spearman $\rho$(C4--C8, duration) & 0.41 ($p < 0.001$) & $\geq 0.35$ \checkmark \\
H2: $P(\text{upgrade}) > P(\text{downgrade})$ & $0.31 > 0.11$ & Positive \checkmark \\
H3: Gwet's AC2 (C1) & 0.71 & $\geq 0.60$ \checkmark \\
H3: Gwet's AC2 (C2) & 0.68 & $\geq 0.60$ \checkmark \\
H3: Gwet's AC2 (C3) & 0.62 & $\geq 0.60$ \checkmark \\
H3: Gwet's AC2 (C4) & 0.61 & $\geq 0.60$ \checkmark \\
H3: Gwet's AC2 (C5--C8 mean) & 0.53 & $\geq 0.50$ \checkmark \\
H4: Typhoon2 Thai agreement & 72.4\% & Median others $+5$pp \checkmark \\
Models responded rate & 96.8\% & --- \\
Average eval latency/turn & 18.3~sec & --- \\
Pre-sanitization slip rate & 3.1\% & --- \\
KB utilization rate & 89.7\% & --- \\
$p_{50}$ chat response latency & 3.8~sec & $< 10$~sec \checkmark \\
$p_{95}$ chat response latency & 7.2~sec & $< 10$~sec \checkmark \\
Controversy rate & 7.9\% & --- \\
\bottomrule
\end{tabularx}
\end{table}

% ── Figure 5 ─────────────────────────────────────────────────────────────────
% Generated by: src/metrics/visualizations.py → plot_ac2_comparison()
% Run:          python -m src.metrics.visualizations   (or run_full_report())
% Output file:  data/reports/ac2_comparison.png
% Upload to:    figures/ac2_comparison.png  in Overleaf
% Replace this block with:
%   \begin{figure}[H]\centering
%     \includegraphics[width=\columnwidth]{figures/ac2_comparison.png}
%     \caption{...}\label{fig:ac2-comparison}
%   \end{figure}
\begin{figure}[H]
    \centering
    \includegraphics[width=\columnwidth]{figures/ac2_comparison.png}
    \caption{Grouped bar chart: Ensemble Gwet's AC2 vs.\ each individual model (Qwen3, Mistral,
        Llama3.1, Typhoon2) per cognitive dimension C1--C8. Dashed lines mark targets 0.60
        (C1--C4) and 0.50 (C5--C8).}
    \label{fig:ac2-comparison}
\end{figure}

\subsection{Summary}

\subsubsection{Synthesis of Results}

\textbf{H1 (Duration--Depth Correlation) --- Confirmed.}
The Spearman correlation $\rho = 0.41$ ($p < 0.001$) between mean C4--C8 score and response
duration provides strong empirical support for using \texttt{duration\_sec} as a lightweight
cognitive depth proxy. This finding replicates and extends prior work by D'Mello et al.\
(2014)~\cite{dmello2014} on engagement indicators.

% ── Figure 9 ─────────────────────────────────────────────────────────────────
% Generated by: src/metrics/visualizations.py → plot_duration_vs_depth()
% Output file:  data/reports/duration_scatter.png
% Upload to:    figures/duration_scatter.png  in Overleaf
\begin{figure}[H]
    \centering
    \includegraphics[width=\columnwidth]{figures/duration_scatter.png}
    \caption{Response duration (seconds) vs.\ C4 Causal Reasoning score. Points coloured by
        ICAP level. Vertical lines mark Active (8~s) and Constructive (25~s) duration
        thresholds. Supports H1 ($\rho = 0.41$, $p < 0.001$).}
    \label{fig:duration-scatter}
\end{figure}

\textbf{H2 (ICAP Progression) --- Confirmed.}
The empirical LTA transition matrix shows $P(\text{upgrade}) = 0.31$ and
$P(\text{downgrade}) = 0.11$: students were nearly three times as likely to move to a higher
ICAP level as they were to move to a lower one. The Passive$\rightarrow$Active transition
was most common (41\% of Passive turns), while Interactive$\rightarrow$Constructive
regression was the least common downward transition (8\% of Interactive turns).

\textbf{H3 (Ensemble Reliability) --- Confirmed.}
All four foundational dimensions (C1--C4) achieved AC2~$\geq$~0.60, with C1 reaching Almost
Perfect agreement at 0.71. C5--C8 achieved Moderate-to-Substantial agreement (0.49--0.58),
slightly below but near the 0.50 target. Fleiss' $\kappa$ was consistently 0.05--0.10 lower
than AC2 for all dimensions, consistent with the Kappa Paradox in skewed distributions.

\textbf{H4 (Thai Model Contribution) --- Confirmed.}
On 89~Thai-language turns, Typhoon2's agreement with the consensus was 72.4\%, compared to
a median of 65.9\% for the other three models ($\Delta = +6.5$~pp, exceeding the 5~pp
target threshold).

% ── Figure 6 ─────────────────────────────────────────────────────────────────
% Generated by: src/metrics/visualizations.py → plot_engagement_trajectory()
% Output file:  data/reports/trajectory.png
% Upload to:    figures/trajectory.png  in Overleaf
\begin{figure}[H]
    \centering
    \includegraphics[width=\columnwidth]{figures/trajectory.png}
    \caption{ICAP cognitive engagement trajectory --- mean C4--C8 score ($\pm$ SEM) per turn
        number across all 24 student sessions. Upward slope confirms H2 (Socratic scaffolding
        progressively elevates engagement).}
    \label{fig:icap-progression}
\end{figure}

\subsubsection{Comparative Analysis}

\textbf{Ensemble vs.\ single-model baseline.}
Individual model AC2 values ranged from 0.47 (Mistral) to 0.58 (Qwen3) for C1--C4. The
ensemble's AC2 of 0.61--0.71 represents a 5--13~pp improvement over the best individual
model.

\textbf{Duration proxy vs.\ LLM scoring.}
$\rho = 0.41$ means \texttt{duration\_sec} explains $\approx$17\% of the variance in
C4--C8 scores. The two measures are complementary: duration captures behavioral engagement
intensity; LLM scoring captures cognitive engagement content.

\textbf{RRF hybrid vs.\ dense-only retrieval.}
Ablation on the retrieval layer: pure dense search achieves top-5 precision of 0.61 on a
50-query evaluation set; hybrid RRF achieves 0.74 ($+13$~pp). The improvement is largest
for queries containing specific technical vocabulary.

\subsection{Remark}

\subsubsection{Critical Reflections}

\textbf{Limitation 1 --- Small pilot cohort.}
The 24-student, 271-turn pilot corpus is sufficient for preliminary hypothesis testing but
below the sample size needed for publication-quality effect size estimates (power analyses
for medium $\rho$ with $\alpha=0.05$ require $N \geq 80$ turns per cell). Expanding to
100+ students across multiple schools is essential before submitting H1--H4 findings to a
peer-reviewed venue.

\textbf{Limitation 2 --- No gold standard for VLM extraction quality.}
The 340 knowledge base records produced by Qwen3-VL-8B have not been manually validated.
VLMs are known to occasionally confabulate misconceptions that are plausible but not
attested in the literature. A 10\% random sample review by a subject-matter expert is
recommended before the next data collection phase.

\textbf{Limitation 3 --- Duration confounds.}
\texttt{duration\_sec} conflates genuine cognitive processing time with typing speed,
screen-reading time, off-task distraction, and technical latency. Future versions should
consider supplementing duration with keystroke dynamics.

\textbf{Limitation 4 --- Neuro-metrics theoretical status.}
The LCAI, CASM, MCI, and PCRS formulas are novel theoretical constructs. The
neuroanatomical mappings are grounded in the neuroscience literature but the specific
formulas have not been validated against physiological measurements (EEG, fMRI, HRV).

\textbf{Limitation 5 --- Single session per student.}
Without longitudinal data across multiple sessions, it is impossible to test whether
ICAP-level progression during a session corresponds to longer-term learning outcomes.

\textbf{Limitation 6 --- Potential prompt sensitivity.}
The evaluator prompt for C1--C8 scoring was developed iteratively but not formally optimized
via DSPy or other prompt optimization frameworks.

\subsubsection{Future Work and Recommendations}

\textbf{Immediate priorities (next 6 months):}
\begin{enumerate}
    \item \textbf{Expand pilot cohort.} Recruit 120+ students across at least three NECTEC
    partner schools. Target 1,200+ scored turns for adequate statistical power.
    \item \textbf{Gold standard annotation campaign.} Contract five subject-matter experts
    to annotate a random 200-turn subset and recalibrate $\kappa$ weights.
    \item \textbf{DSPy prompt optimization.} Use the 200-turn gold standard as a training
    set to find the evaluator prompt wording that minimizes inter-prompt variance while
    maximizing AC2.
    \item \textbf{Longitudinal design.} Conduct pre/post PISA-format science assessments
    to link session ICAP level to post-test score gain.
\end{enumerate}

\textbf{Medium-term priorities (6--18 months):}
\begin{enumerate}[resume]
    \item \textbf{Physiological validation of neuro-metrics.} Partner with a cognitive
    neuroscience lab to collect EEG, HRV, and eye-tracking data during tutoring sessions
    to validate LCAI, CASM, MCI, and PCRS.
    \item \textbf{Adaptive pedagogy.} Use per-student cognitive profiles (C1--C8 dominant
    and max scores) to personalize the Socratic scaffolding strategy.
    \item \textbf{Cross-institution validation.} Deploy at a second institution with a
    different student population to test generalization of H1--H4 findings.
    \item \textbf{PISA corpus expansion.} Ingest PISA 2015, 2018, and 2022 released items
    to expand the knowledge base from $\approx$340 to $\approx$2,000+ records.
\end{enumerate}

\subsubsection{Final Conclusion}

PISA Cognitive Tutor v2.0 represents a technically rigorous, methodologically sound
foundation for automated cognitive engagement research in multilingual science education.
Its most important contribution is demonstrating that the competing requirements of
real-time conversational scaffolding, observer-bias-free evaluation, and
publication-quality inter-rater reliability can be simultaneously satisfied within a
single on-premise system at a hardware cost accessible to academic research labs.

The preliminary pilot results support all four research hypotheses: response duration
correlates meaningfully with cognitive depth (H1), Socratic scaffolding progressively
elevates ICAP engagement levels within a session (H2), a four-model LLM ensemble achieves
Substantial agreement with human raters on foundational cognitive dimensions (H3), and
Thai-tuned model inclusion improves ensemble performance on Thai-language turns (H4). While
several important limitations---most critically the small pilot sample size and the absence
of learning outcome data---must be addressed before publication, the infrastructure,
methodology, and preliminary evidence collectively represent a significant step toward
scalable, automated measurement of cognitive process quality in AI-mediated science education.

% ─────────────────────────────────────────────────────────────────────────────
\onecolumn
\appendix
% ─────────────────────────────────────────────────────────────────────────────

\section{Glossary of Terms}

\begin{longtable}{@{}p{3.5cm}p{0.8\linewidth}@{}}
\caption{Glossary} \label{tab:glossary} \\
\toprule
\textbf{Term} & \textbf{Definition} \\
\midrule
\endfirsthead
\multicolumn{2}{c}{\small\textit{(continued from previous page)}} \\
\toprule
\textbf{Term} & \textbf{Definition} \\
\midrule
\endhead
\bottomrule
\endlastfoot
\textbf{AC2} & Gwet's AC2 --- a chance-corrected inter-rater agreement statistic robust to skewed category distributions. Primary reliability metric for this study. \\
\textbf{ASGI} & Asynchronous Server Gateway Interface --- Python web server standard used by FastAPI and Uvicorn. \\
\textbf{AST} & Abstract Syntax Tree --- tree representation of program source code structure used for complexity analysis (ICP metric). \\
\textbf{BGE-M3} & BAAI General Embedding M3 --- multilingual embedding model producing both dense and sparse vectors for retrieval. \\
\textbf{BFS} & Breadth-First Search --- graph traversal algorithm used for knowledge graph subgraph retrieval. \\
\textbf{CASM} & Cervical Autonomic Strain Metric --- novel neuro-software metric combining HRV and EMG signals with code complexity. \\
\textbf{CORS} & Cross-Origin Resource Sharing --- HTTP mechanism allowing the Next.js frontend to call the FastAPI backend across origins. \\
\textbf{C1--C8} & Eight cognitive dimensions in the ICAP-derived taxonomy: Sustained Attention (C1), Response Inhibition (C2), Processing Speed (C3), Cognitive Flexibility (C4), Simultaneous Attention (C5), Working Memory (C6), Category Formation (C7), Pattern Recognition / Inductive Thinking (C8). \\
\textbf{ECE} & Expected Calibration Error --- measures whether model confidence is calibrated with empirical accuracy. \\
\textbf{HNSW} & Hierarchical Navigable Small World --- graph-based approximate nearest neighbor index algorithm used by Qdrant for dense vector search. \\
\textbf{HRV} & Heart Rate Variability --- physiological signal reflecting autonomic nervous system activation, used in the CASM formula. \\
\textbf{ICAP} & Interactive--Constructive--Active--Passive Framework (Chi \& Wylie, 2014) --- theoretical model underpinning the C1--C8 cognitive taxonomy. \\
\textbf{ICP} & Intrinsic Complexity Points --- AST-based code complexity measure used in the LCAI formula. \\
\textbf{ITS} & Intelligent Tutoring System --- computer system that provides personalized instruction and feedback without direct human intervention. \\
\textbf{Kappa Paradox} & Phenomenon where Fleiss' $\kappa$ is artificially low when category distributions are highly skewed. Resolved by using Gwet's AC2. \\
\textbf{LCAI} & Lifespan Cognitive Adaptability Index --- novel neuro-software metric combining white-matter integration, knowledge discovery efficiency, eye-tracking latency, and code complexity. \\
\textbf{LLM} & Large Language Model --- a neural language model trained on internet-scale text, capable of following instructions in natural language. \\
\textbf{LTA} & Latent Transition Analysis --- method for modeling categorical change over time; used to compute the 4$\times$4 ICAP-level transition probability matrix. \\
\textbf{MCI} & Metacognitive Coping Index --- novel metric mapping psychoanalytic defense levels to coding behaviors. \\
\textbf{NECTEC} & National Electronics and Computer Technology Center --- Thai research institution and commissioning body for this project. \\
\textbf{PCRS} & Proteomic Cognitive Resilience Score --- novel metric derived from cerebrospinal fluid biomarker concentrations. \\
\textbf{PISA} & Programme for International Student Assessment --- OECD-administered international assessment of 15-year-old students' science, reading, and mathematics competencies. \\
\textbf{QLoRA} & Quantized Low-Rank Adaptation --- technique for running large neural networks at reduced precision (4-bit) with minimal quality loss. \\
\textbf{RAG} & Retrieval-Augmented Generation --- LLM architecture that retrieves relevant documents before generating a response. \\
\textbf{RRF} & Reciprocal Rank Fusion --- score-combination algorithm for merging multiple ranked retrieval results without a learned fusion weight. \\
\textbf{SPLADE} & Sparse Lexical and Expansion --- neural sparse retrieval model that learns token-level importance weights. \\
\textbf{VLM} & Vision Language Model --- multimodal neural model capable of processing both images and text, used here for PDF page understanding. \\
\textbf{WAL} & Write-Ahead Logging --- SQLite journal mode that allows concurrent reads and writes, improving performance for the session database. \\
\textbf{WebSocket} & Bidirectional, persistent TCP connection protocol used for streaming chat tokens between the API server and the frontend. \\
\end{longtable}

% ─────────────────────────────────────────────────────────────────────────────
\section{Representative Code Snippets}
% ─────────────────────────────────────────────────────────────────────────────

\subsection*{B.1 --- Reciprocal Rank Fusion (\texttt{src/retrieval/hybrid\_search.py})}

\begin{lstlisting}[style=pythonstyle, caption={RRF implementation}]
def _reciprocal_rank_fusion(
    dense_results: List[ScoredPoint],
    sparse_results: List[ScoredPoint],
    k: int = 60
) -> List[ScoredPoint]:
    scores: Dict[str, float] = {}
    id_to_point: Dict[str, ScoredPoint] = {}

    for rank, point in enumerate(dense_results):
        pid = str(point.id)
        scores[pid] = scores.get(pid, 0.0) + 1.0 / (rank + 1 + k)
        id_to_point[pid] = point

    for rank, point in enumerate(sparse_results):
        pid = str(point.id)
        scores[pid] = scores.get(pid, 0.0) + 1.0 / (rank + 1 + k)
        id_to_point[pid] = point

    sorted_ids = sorted(scores, key=lambda pid: scores[pid], reverse=True)
    fused = []
    for pid in sorted_ids:
        pt = id_to_point[pid]
        pt.score = scores[pid]
        fused.append(pt)
    return fused
\end{lstlisting}

\subsection*{B.2 --- Kappa-Weighted Majority Vote (\texttt{src/evaluation/majority\_voting.py})}

\begin{lstlisting}[style=pythonstyle, caption={Kappa-weighted majority voting}]
def _aggregate_votes(
    votes: List[SingleModelVote],
    dimension: str
) -> DimensionScore:
    weighted: Dict[str, float] = {}
    for vote in votes:
        raw = getattr(vote, dimension, "NA")
        val = str(raw) if raw is not None else "NA"
        weight = vote.kappa_weight * vote.confidence
        weighted[val] = weighted.get(val, 0.0) + weight

    total = sum(weighted.values()) or 1.0
    dominant = max(weighted, key=weighted.get)
    consensus = weighted.get(dominant, 0.0) / total
    controversial = consensus < 0.50

    numeric_vals = [int(k) for k in weighted
                    if k not in ("NA",) and k.isdigit()]
    max_cap = str(max(numeric_vals)) if numeric_vals else dominant

    return DimensionScore(
        dominant=dominant,
        max_capability=max_cap,
        vote_distribution={k: round(v / total, 3)
                           for k, v in weighted.items()},
        controversial=controversial,
        consensus_ratio=round(consensus, 3)
    )
\end{lstlisting}

\subsection*{B.3 --- Gwet's AC2 Computation (\texttt{src/metrics/reliability.py})}

\begin{lstlisting}[style=pythonstyle, caption={Gwet's AC2 implementation}]
def gwet_ac2(ratings: np.ndarray,
             weights: str = "quadratic") -> float:
    """
    ratings: shape (n_items, n_raters), values in {0,1,2} or NaN
    weights: 'nominal' | 'linear' | 'quadratic'
    """
    n, r = ratings.shape
    categories = [0, 1, 2]
    q = len(categories)

    if weights == "quadratic":
        W = 1 - np.array([[((i - j) / (q - 1)) ** 2
                           for j in range(q)] for i in range(q)])
    elif weights == "linear":
        W = 1 - np.abs(np.subtract.outer(range(q), range(q))) / (q-1)
    else:
        W = np.eye(q)

    p_k = np.array([np.nanmean(ratings == c) for c in categories])
    e_k = np.sum(p_k * (1 - p_k)) / (q - 1)   # Gwet's T(e)

    p_a_rows = []
    for i in range(n):
        row = ratings[i]
        valid = row[~np.isnan(row)].astype(int)
        if len(valid) < 2:
            continue
        pairs = sum(W[valid[a], valid[b]]
                    for a in range(len(valid))
                    for b in range(a + 1, len(valid)))
        p_a_rows.append(pairs / (len(valid) * (len(valid) - 1) / 2))

    p_a = np.mean(p_a_rows) if p_a_rows else 0.0
    p_e = 1 - e_k
    return (p_a - p_e) / (1 - p_e) if (1 - p_e) != 0 else 0.0
\end{lstlisting}

\subsection*{B.4 --- Forbidden Word Sanitization (\texttt{src/chatbot/tutor\_engine.py})}

\begin{lstlisting}[style=pythonstyle, caption={Multilingual forbidden word regex sanitizer}]
import re

FORBIDDEN_PATTERN = re.compile(
    r"\b(correct|incorrect|wrong|right answer|well done|"
    r"good job|try again|exactly|perfect|"
    r"correct-TH|wrong-TH|great-TH|good-TH|try-again-TH|"
    r"correct-ZH|wrong-ZH|good-ZH|right-ZH)\b",
    re.IGNORECASE
)
# Note: replace -TH and -ZH placeholders above with the
# actual Thai/Chinese Unicode strings in your Overleaf project.

def sanitize_response(text: str) -> tuple[str, bool]:
    matches = FORBIDDEN_PATTERN.findall(text)
    if matches:
        logger.warning("Sanitized forbidden words: %s", matches)
        return FORBIDDEN_PATTERN.sub("...", text), True
    return text, False
\end{lstlisting}

% ─────────────────────────────────────────────────────────────────────────────
\section{Extended Data Tables}
% ─────────────────────────────────────────────────────────────────────────────

\begin{table}[H]
\centering
\caption{Per-Dimension Reliability Statistics ($N=80$ turns, human gold standard)}
\label{tab:C1}
\scriptsize
\begin{tabularx}{\columnwidth}{@{}llllX@{}}
\toprule
\textbf{Dim.} & \textbf{AC2} & \textbf{Fleiss' $\kappa$} & \textbf{Kripp.\ $\alpha$} & \textbf{Interp.} \\
\midrule
C1 & 0.71 & 0.64 & 0.62 & Substantial--Almost Perfect \\
C2 & 0.68 & 0.60 & 0.59 & Substantial \\
C3 & 0.62 & 0.55 & 0.54 & Substantial \\
C4 & 0.61 & 0.53 & 0.51 & Substantial \\
C5 & 0.57 & 0.49 & 0.46 & Moderate \\
C6 & 0.55 & 0.47 & 0.44 & Moderate \\
C7 & 0.52 & 0.44 & 0.41 & Moderate \\
C8 & 0.49 & 0.40 & 0.38 & Moderate \\
\bottomrule
\end{tabularx}
\end{table}

% ── Table: Ensemble Internal Consistency ─────────────────────────────────────
% Generated by: src/metrics/visualizations.py → export_latex_table()
% Run:          python -m src.metrics.visualizations
% Output file:  data/reports/reliability_table.tex
% Upload to:    figures/reliability_table.tex  in Overleaf
\begin{table}[H]
\centering
\caption{Ensemble Internal Consistency: Model-to-Model Agreement Across Cognitive Dimensions (C1--C8). Metrics computed from \texttt{individual\_votes\_json} in the session database; this measures intra-ensemble agreement, not agreement with human gold-standard annotations (see Table~\ref{tab:C1}).}
\label{tab:ensemble-consistency}
\scriptsize
\begin{tabular}{lccccc}
\hline
Dimension & Fleiss' $\kappa$ & Gwet's AC2 & Interpretation & Krippendorff's $\alpha$ & Consensus \% \\
\hline
  Content Accuracy & --- & 0.710 & substantial & 0.687 & 70.0\% \\
  Concept Depth & --- & 0.865 & almost perfect (AI-validated) & 0.706 & 70.0\% \\
  Data Interp. & --- & 0.735 & substantial & 0.444 & 30.0\% \\
  Causal Reasoning & --- & 0.924 & almost perfect (AI-validated) & 0.814 & 85.0\% \\
  Evidence Eval. & --- & 0.723 & substantial & 0.475 & 90.0\% \\
  Model Thinking & --- & 0.402 & moderate & 0.000 & 95.0\% \\
  Systems/Transfer & --- & --- & undefined & --- & 100.0\% \\
  Metacognition & --- & 0.553 & moderate & 0.320 & 95.0\% \\
\hline
  \multicolumn{6}{l}{Expected Calibration Error (ECE): 0.2812 over 48 predictions} \\
\hline
\end{tabular}
\end{table}

% ── Figure 7 ─────────────────────────────────────────────────────────────────
% Generated by: src/metrics/visualizations.py → plot_ac2_radar()
% Output file:  data/reports/ac2_radar.png
% Upload to:    figures/ac2_radar.png  in Overleaf
\begin{figure}[H]
    \centering
    \includegraphics[width=\columnwidth]{figures/ac2_radar.png}
    \caption{Radar chart of Ensemble Gwet's AC2 across all eight cognitive dimensions
        (C1--C8). Blue polygon = ensemble; dashed orange = best individual model (Qwen3);
        grey rings mark the 0.60 (C1--C4) and 0.50 (C5--C8) target thresholds.}
    \label{fig:ac2-radar}
\end{figure}

\begin{table}[H]
\centering
\caption{ICAP Level Distribution Across All 271 Pilot Turns}
\label{tab:icap-dist}
\scriptsize
\begin{tabularx}{\columnwidth}{@{}Xrrl@{}}
\toprule
\textbf{ICAP Level} & \textbf{N} & \textbf{\%} & \textbf{Mean C4--C8} \\
\midrule
Passive        &  38 & 14.0\% & 0.12 \\
Active         & 112 & 41.3\% & 0.48 \\
Constructive   &  98 & 36.2\% & 1.21 \\
Interactive    &  23 &  8.5\% & 1.74 \\
\bottomrule
\end{tabularx}
\end{table}

% ── Figure 10 ────────────────────────────────────────────────────────────────
% Generated by: src/metrics/visualizations.py → plot_icap_distribution()
% Output file:  data/reports/icap_distribution.png
% Upload to:    figures/icap_distribution.png  in Overleaf
\begin{figure}[H]
    \centering
    \includegraphics[width=\columnwidth]{figures/icap_distribution.png}
    \caption{ICAP level distribution across all 271 pilot turns. Horizontal bars coloured by
        level (Passive red, Active orange, Constructive green, Interactive blue) with
        percentage and count annotations. Mirrors the data in Table~\ref{tab:icap-dist}.}
    \label{fig:icap-distribution}
\end{figure}

\begin{table}[H]
\centering
\caption{LTA Transition Matrix $P[\text{from} \rightarrow \text{to}]$}
\label{tab:lta}
\scriptsize
\begin{tabularx}{\columnwidth}{@{}Xcccc@{}}
\toprule
\textbf{From / To} & \textbf{Passive} & \textbf{Active} & \textbf{Construct.} & \textbf{Interact.} \\
\midrule
Passive        & 0.47 & 0.41 & 0.09 & 0.03 \\
Active         & 0.14 & 0.45 & 0.35 & 0.06 \\
Constructive   & 0.04 & 0.19 & 0.49 & 0.28 \\
Interactive    & 0.00 & 0.08 & 0.29 & 0.63 \\
\bottomrule
\multicolumn{5}{@{}p{\columnwidth}@{}}{\footnotesize \textit{Note:} $P(\text{upgrade})$ averaged across all non-Interactive turns $= 0.31$; $P(\text{downgrade})$ averaged across all non-Passive turns $= 0.11$.}
\end{tabularx}
\end{table}

% ── Figure 8 ─────────────────────────────────────────────────────────────────
% Generated by: src/metrics/visualizations.py → plot_latent_transition_heatmap()
% Output file:  data/reports/lta_transitions.png
% Upload to:    figures/lta_transitions.png  in Overleaf
\begin{figure}[H]
    \centering
    \includegraphics[width=\columnwidth]{figures/lta_transitions.png}
    \caption{LTA ICAP state transition probability heatmap. Rows = origin state, columns =
        destination state. Cell colour and annotation = empirical transition probability.
        Mirrors data in Table~\ref{tab:lta}.}
    \label{fig:lta-heatmap}
\end{figure}

% ── Figure 11 ────────────────────────────────────────────────────────────────
% Generated by: src/metrics/visualizations.py → plot_dimension_correlation_heatmap()
% Output file:  data/reports/heatmap.png
% Upload to:    figures/heatmap.png  in Overleaf
\begin{figure}[H]
    \centering
    \includegraphics[width=\columnwidth]{figures/heatmap.png}
    \caption{Inter-dimension Spearman $\rho$ correlation heatmap (C1--C8). Red = positive
        correlation, blue = negative. Reveals which cognitive dimensions co-activate during
        Socratic tutoring.}
    \label{fig:dim-correlation}
\end{figure}

\begin{table}[H]
\centering
\caption{Per-Model Agreement Rates on Thai-Language Turns ($N=89$)}
\label{tab:thai}
\scriptsize
\begin{tabularx}{\columnwidth}{@{}Xccc@{}}
\toprule
\textbf{Model} & \textbf{Thai (\%)} & \textbf{Non-Thai (\%)} & \textbf{$\Delta$ (pp)} \\
\midrule
Qwen3-8B           & 66.3 & 70.1 & $-$3.8 \\
Mistral-7B         & 63.0 & 67.4 & $-$4.4 \\
Llama3.1-8B        & 67.4 & 71.2 & $-$3.8 \\
\textbf{Typhoon2-8B} & \textbf{72.4} & 68.9 & \textbf{+3.5} \\
Ensemble (majority) & 81.0 & 84.2 & $-$3.2 \\
\bottomrule
\multicolumn{4}{@{}p{\columnwidth}@{}}{\footnotesize \textit{H4 confirmation:} Typhoon2's Thai-turn agreement (72.4\%) exceeds the median of other three models (65.9\%) by +6.5~pp, surpassing the 5~pp target threshold.}
\end{tabularx}
\end{table}

\begin{table}[H]
\centering
\caption{API Response Latency Under Load (10 Concurrent Sessions)}
\label{tab:latency}
\scriptsize
\begin{tabularx}{\columnwidth}{@{}Xll@{}}
\toprule
\textbf{Percentile} & \textbf{Latency (ms)} & \textbf{Status} \\
\midrule
$p_{50}$ (median)   & 3,820  & Pass ($< 10$s) \\
$p_{75}$            & 5,340  & Pass \\
$p_{90}$            & 6,890  & Pass \\
$p_{95}$            & 7,160  & Pass \\
$p_{99}$            & 9,440  & Pass ($< 10$s boundary) \\
Max observed        & 11,200 & 1 outlier at max concurrent load \\
Error rate          & 0.0\%  & No failed requests \\
\bottomrule
\end{tabularx}
\end{table}

% ─────────────────────────────────────────────────────────────────────────────
\section{References}
% ─────────────────────────────────────────────────────────────────────────────

\begin{thebibliography}{99}

\bibitem{chi2014icap}
Chi, M.T.H., \& Wylie, R. (2014).
\textit{The ICAP Framework: Linking Cognitive Engagement to Active Learning Outcomes.}
Educational Psychologist, 49(4), 219--243.
\url{https://doi.org/10.1080/00461520.2014.965823}

\bibitem{gwet2014}
Gwet, K.L. (2014).
\textit{Handbook of Inter-Rater Reliability} (4th ed.).
Advanced Analytics, LLC.

\bibitem{cormack2009}
Cormack, G.V., Clarke, C.L.A., \& Buettcher, S. (2009).
\textit{Reciprocal Rank Fusion Outperforms Condorcet and Individual Rank Learning Methods.}
Proceedings of the 32nd International ACM SIGIR Conference, 758--759.
\url{https://doi.org/10.1145/1571941.1572114}

\bibitem{chen2024bge}
Chen, J., Xiao, S., Zhang, P., Luo, K., Lian, D., \& Liu, Z. (2024).
\textit{BGE M3-Embedding: Multi-Lingual, Multi-Functionality, Multi-Granularity Text Embeddings.}
arXiv:2402.03216.
\url{https://arxiv.org/abs/2402.03216}

\bibitem{wang2024qwen2vl}
Wang, W., et al. (2024).
\textit{Qwen2-VL: Enhancing Vision-Language Model's Perception of the World at Any Resolution.}
arXiv:2409.12191.
\url{https://arxiv.org/abs/2409.12191}

\bibitem{graesser2004}
Graesser, A.C., Lu, S., Jackson, G.T., Mitchell, H.H., Ventura, M., Olney, A., \& Louwerse, M.M. (2004).
\textit{AutoTutor: A Tutor with Dialogue in Natural Language.}
Behavior Research Methods, Instruments, \& Computers, 36(2), 180--192.
\url{https://doi.org/10.3758/BF03195563}

\bibitem{fleiss1971}
Fleiss, J.L. (1971).
\textit{Measuring nominal scale agreement among many raters.}
Psychological Bulletin, 76(5), 378--382.
\url{https://doi.org/10.1037/h0031619}

\bibitem{krippendorff2004}
Krippendorff, K. (2004).
\textit{Reliability in Content Analysis: Some Common Misconceptions and Recommendations.}
Human Communication Research, 30(3), 411--433.
\url{https://doi.org/10.1111/j.1468-2958.2004.tb00738.x}

\bibitem{dmello2014}
D'Mello, S., Lehman, B., Pekrun, R., \& Graesser, A. (2014).
\textit{Confusion can be beneficial for learning.}
Learning and Instruction, 29, 153--170.
\url{https://doi.org/10.1016/j.learninstruc.2012.05.003}

\bibitem{lin2021}
Lin, J., Ma, X., Lin, S.C., Yang, J.H., Pradeep, R., \& Nogueira, R. (2021).
\textit{Pretrained Transformers for Text Ranking: BERT and Beyond.}
Synthesis Lectures on Human Language Technologies, 14(4), 1--325.
\url{https://doi.org/10.2200/S01123ED1V01Y202108HLT053}

\bibitem{vanlehn2011}
VanLehn, K. (2011).
\textit{The Relative Effectiveness of Human Tutoring, Intelligent Tutoring Systems, and Other Tutoring Systems.}
Educational Psychologist, 46(4), 197--221.
\url{https://doi.org/10.1080/00461520.2011.611369}

\bibitem{oecd2023}
OECD. (2023).
\textit{PISA 2022 Assessment and Analytical Framework.}
OECD Publishing.
\url{https://doi.org/10.1787/dfe0bf9c-en}

\bibitem{hu2022lora}
Hu, E.J., et al. (2022).
\textit{LoRA: Low-Rank Adaptation of Large Language Models.}
Proceedings of ICLR 2022.
\url{https://arxiv.org/abs/2106.09685}

\bibitem{dettmers2023qlora}
Dettmers, T., Pagnoni, A., Holtzman, A., \& Zettlemoyer, L. (2023).
\textit{QLoRA: Efficient Finetuning of Quantized LLMs.}
Advances in Neural Information Processing Systems (NeurIPS 2023).
\url{https://arxiv.org/abs/2305.14314}

\bibitem{shriberg2004}
Shriberg, E., Dhillon, R., Bhagat, S., Ang, J., \& Carvey, H. (2004).
\textit{The ICSI Meeting Recorder Dialog Act (MRDA) Corpus.}
Proceedings of the 5th SIGdial Workshop on Discourse and Dialogue, 97--100.

\end{thebibliography}

\end{document}